	\begin{UseCase}{CU9}{Gestionar Mensajes}{
		Cuando un colaborador (\hyperlink{jefe}{Líder de análisis} o \hyperlink{analista}{Analista}) está operando sobre un \hyperlink{proyectoEntidad}{Proyecto}, tiene la posibilidad de gestionar los \hyperlink{tElemento}{elementos} que serán referenciados al \hyperlink{casoUso}{Caso de Uso}. Uno de estos \hyperlink{tElemento}{elementos} son los Mensajes.\\
		
		Un \hyperlink{MSGEntidad}{Mensaje} es un texto que comprende la mínima unidad de comunicación entre el usuario y el sistema, estos se encargan a través del sistema de informar a los actores sobre errores, avisos, informes y notificaciones.\\
		
		Este caso de uso permite al colaborador visualizar en una tabla, el registro de los mensajes que forman parte del proyecto, así como solicitar el registro de un mensaje nuevo, modificar alguno existente (teniendo la posibilidad de actualizar algún dato de registro), eliminarlo (en caso de que el mensaje por algún motivo ya no forme parte del proyecto) o simplemente consultar su información (nombre, descripción y redacción del mensaje) \\
		Las acciones disponibles para cada \hyperlink{MSGEntidad}{Mensaje} dependerán del estado en el que se encuentre el caso de uso donde es referenciado. 
	}

	\UCitem{Actor}{\hyperlink{jefe}{Líder de análisis}, \hyperlink{analista}{Analista}}
	\UCitem{Propósito}{Proporcionar al actor un mecanismo para llevar el control de los mensajes de un proyecto.}
	\UCitem{Entradas}{Ninguna}
	\UCitem{Salidas}{\begin{itemize}
			\item \cdtRef{proyectoEntidad:claveProyecto}{Clave del proyecto}: Lo obtiene el sistema.
			\item \cdtRef{proyectoEntidad:nombreProyecto}{Nombre del proyecto}: Lo obtiene el sistema.
			\item \cdtRef{MSGEntidad}{Mensaje}: Tabla que muestra \cdtRef{MSGEntidad:nombreMSG}{nombre} de todos los mensajes registrados de un proyecto.
			\item \cdtIdRef{MSG2}{No existe información}: Se muestra en la pantalla \IUref{IU10A}{Gestionar Mensajes} cuando no existen mensajes registradas.
	\end{itemize}}
	
	\UCitem{Precondiciones}{Que exista al menos un proyecto registrado.}
	\UCitem{Postcondiciones}{Ninguna}
	\UCitem{Errores}{Ninguno}
	\UCitem{Tipo}{Primario}
\end{UseCase}
%--------------------------------------------------------------------------------------------------------------------------------------------------------------------------
\begin{UCtrayectoria}
	\UCpaso[\UCactor] Solicita gestionar los mensajes seleccionando la opción ''Mensajes'' del menú \IUref{MN3}{Menú de Proyecto}.
	\UCpaso[\UCsist] Obtiene la información de los mensajes registrados del proyecto seleccionado. \hyperlink{CU9:TAA}{[Trayectoria A]}
	\UCpaso[\UCsist] Ordena los mensajes alfabéticamente basándose en la clave de los mismos.
	\UCpaso[\UCsist] Muestra la información de los mensajes en la pantalla \IUref{IU10}{Gestionar Mensajes} y las operaciones disponibles de acuerdo a la regla de negocio \BRref{RN15}{Operaciones disponibles}.
	\UCpaso[\UCactor] Gestiona los mensajes a través de los botones: \IUbutton{Registrar}, \editar y \eliminar. \label{CU9-P4}
\end{UCtrayectoria}		
%--------------------------------------------------------------------------------------------------------------------------------
\hypertarget{CU9:TAA}{\textbf{Trayectoria alternativa A}}\\
\noindent \textbf{Condición:} No existen registros de Mensajes.
\begin{enumerate}
	\UCpaso[\UCsist] Muestra el mensaje \cdtIdRef{MSG2}{No existe información} en la pantalla \IUref{IU10A}{Gestionar Mensajes} para indicar que no hay registros de mensajes para mostrar. \label{CU9-TA1}
	\UCpaso[\UCactor] Gestiona los mensajes a través del botón: \IUbutton{Registrar}. 
	\item[- -] - - {\em {Fin del caso de uso}}.%
\end{enumerate}

%--------------------------------------

\subsubsection{Puntos de extensión}

\UCExtenssionPoint{El actor requiere registrar un mensaje.}{Presionando el botón \IUbutton{Registrar} del paso \ref{CU9-P4} de la trayectoria principal o del paso \ref{CU9-TA1} de la trayectoria alternativa A.}{\UCref{CU9.1}{Registrar Mensaje}}
\UCExtenssionPoint{El actor requiere modificar un mensaje.}{Presionando el icono \editar del paso \ref{CU9-P4} de la trayectoria principal.}{\UCref{CU9.2}{Modificar Mensaje}}
\UCExtenssionPoint{El actor requiere eliminar un mensae.}{Presionando el icono \eliminar del paso \ref{CU9-P4} de la trayectoria principal.}{\UCref{CU9.3}{Eliminar Mensaje}}
\UCExtenssionPoint{El actor requiere consultar un mensaje.}{Presionando el icono \raisebox{-1mm}{\includegraphics[height=11pt]{images/Iconos/consultar}} del paso \ref{CU9-P4} de la trayectoria principal.}{\UCref{CU9.4}{Consultar Mensaje}}